\section{Design of Hetero-KVCache-Optimizer}
\label{sec:design}

The core philosophy of \texttt{Hetero-KVCache-Optimizer} is to decouple the \textit{logical computation graph} from the \textit{physical memory layout}. This allows the engine to perform full-sequence attention during prefill while maintaining a strictly bounded memory footprint during decoding, even with visual tokens from multiple video frames.

\subsection{Transient Cache Interception}
To bypass the "Flash-Fallback" problem, we design a transient interception layer[cite: 3]. Our system treats the prefill KV tensors as \textit{short-lived (transient)} activations[cite: 4].

% --- Figure 1 Placeholder: System Architecture ---
\begin{figure}[htbp]
  \centering
  % Replace filename with actual figure name after generation
  \includegraphics[width=\linewidth]{figures/system_architecture.png}
  \caption{Overall Architecture of Hetero-KVCache-Optimizer. The Transient Interceptor hooks the prefill stream to extract persistent KV tokens while maintaining FlashAttention compatibility[cite: 5].}
  \label{fig:arch}
\end{figure}

\begin{algorithm}[h]
\caption{Transient Interception \& Pruning}
\label{alg:transient}
\begin{algorithmic}[1]
\STATE \textbf{Input:} Full input sequence $X_{1 \dots N}$
\STATE \textbf{Output:} Optimized KV Pool $\mathcal{P}$, Hidden State $H$
\STATE $K_{full}, V_{full} \gets \text{Prefill}(X_{1 \dots N})$ 
\STATE $H \gets \text{FlashAttention}(Q, K_{full}, V_{full})$ 
\STATE // \textit{Extraction and Immediate GC}
\STATE $\mathcal{S} \gets \text{ExtractSinks}(K_{full}, V_{full}, \text{count}=64)$
\STATE $\mathcal{T} \gets \text{ExtractTail}(K_{full}, V_{full}, \text{count}=8192)$
\STATE $\mathcal{P}.\text{Update}(\mathcal{S} \cup \mathcal{T})$ [cite: 5]
\STATE \textbf{Free}($K_{full}, V_{full}$) \COMMENT{Immediate GC to reclaim peak VRAM}
\RETURN $H, \mathcal{P}$
\end{algorithmic}
\end{algorithm}

The massive $K_{full}$ and $V_{full}$ tensors only exist during the forward pass. Once computed, the system "harvests" essential Sinks and Tail tokens, allowing explicit \texttt{del} operations and \texttt{gc.collect()} to reclaim gigabytes of memory immediately.

\subsection{Static Zero-Fragmentation Pool}
To eliminate Python's dynamic allocation overhead, we implement a \textbf{Static Physical Pool}[cite: 8]. We pre-allocate a contiguous VRAM block upon initialization[cite: 9].

% --- Figure 2 Placeholder: Memory Flow ---
\begin{figure}[htbp]
  \centering
  \includegraphics[width=0.85\linewidth]{figures/memory_pool_mechanism.png}
  \caption{Static Memory Pool with In-place Shift Mechanism. The layout remains constant throughout the 70k token process to eliminate external fragmentation[cite: 10, 12].}
  \label{fig:memory}
\end{figure}

Instead of \texttt{torch.cat}, we use an \textbf{In-place Shift} mechanism[cite: 10, 11]. When the tail buffer is full, we perform a pointer-based shift or circular indexing[cite: 12]:
\begin{equation}
\text{Index}_{phys} = (\text{Index}_{logical} + \text{Offset}) \pmod{L_{max}}
\end{equation}

\subsection{Strict Mask \& RoPE Alignment}
Reducing the KV Cache typically destroys relative distance in RoPE[cite: 13]. We implement a \textbf{Coordinate Re-mapping} strategy to maintain 100\% accuracy[cite: 14]:
\begin{itemize}
    \item \textbf{Attention Mask Synchronization}: Blocks attention to evicted middle tokens while maintaining causality[cite: 15].
    \item \textbf{RoPE Preservation}: Caches original $sin/cos$ values alongside KV pairs to retain absolute positional information.
\end{itemize}